% !TEX root = ../../../tree_notes.tex
% =====================================================================
%  15.4  Rerooting a count / a per-node verdict   sec:c20-counting
%  Exports: def:c20-count-state  thm:c20-count-edge-transfer  lst:c20-count
%           fig:c20-countfig  ex:c20-count-run  prob:c20-count-solved
%           lst:c20-1626e-soln
%  The combine here is a plain SUM (invertible), so the "up-pass" collapses to
%  one subtraction: TOTAL - down[c]. That is the whole parent-side piece. The
%  interesting move is reading the rerooted count on EACH SIDE of EVERY edge and
%  turning it into a yes/no verdict per node.
%  Anchor: CF 1626E Black and White Tree -- a chip hops toward a chosen black
%  vertex (never the same black twice in a row); report for every start vertex
%  whether it can reach a black vertex. Editorial (Educational Round 121) turns
%  this into a per-edge count (blacks on the far side of each cut) plus a
%  multisource reachability. Official n=8 sample gives 0 1 1 1 1 0 1 1.
%  Overlay hand-computed on the running tree T with black = {4, 10}.
% =====================================================================
\section{Rerooting a Count: A Verdict for Every Node}\label{sec:c20-counting}

\begin{previouslybox}
\begin{enumerate}[nosep, leftmargin=1.6em]
  \item In \secref{c20-updown} the parent-side piece of an invertible $\Phi$ was
        one subtraction. For a \emph{count of marked nodes}, what exactly is
        subtracted from what? (the child's subtree count from the total count)
  \item Say the tree has $B$ marked nodes and edge $(v,c)$ is cut, $c$ on the
        child side. How many marks sit on $v$'s side, in one expression?
        ($B - \mathit{down}[c]$, where $\mathit{down}[c]$ is the marks in
        $\subtree(c)$)
  \item Ch.~13's fold gives $\mathit{down}[c]$ for a fixed root. Why is the
        \emph{other} side the only thing a reroot buys you here? (the child side
        is unchanged by re-rooting; the complement is the up-value)
  \item A \emph{verdict} (yes/no) is not a sum. What does the count feed, if not
        a running total? (a per-edge \emph{comparison} that decides a direction)
\end{enumerate}
\end{previouslybox}

\begin{topicintro}
The last three sections rerooted \emph{values} you keep and re-combine: a
distance sum, a pair of component sizes, a signed product. This section reroots
a \emph{count} -- and then does something different with it. The count of marked
nodes in a subtree is a plain sum, so its parent-side piece is a single
subtraction $\mathit{TOTAL}-\mathit{down}[c]$; nothing new there. The new idea is
that the answer you want is not the count at a node but a \term{verdict} read off
the count on \emph{both sides of every edge}: ``is this edge safe to cross?''
``can this node reach a target?'' Those yes/no questions are decided by
\emph{comparing} the two rerooted counts, and the per-node verdicts are then
assembled by a reachability sweep over the directed edges that passed the test.
We build that pattern (a rerooted count feeding a per-edge decision, \defref
{c20-count-state}, \thmref{c20-count-edge-transfer}) and spend it on CF~1626E:
\emph{a chip hops toward a chosen black vertex, never the same black twice in a
row; decide for every start vertex whether some black vertex is reachable}. It is
the count-reroot cousin of the game values that reach their full form in this
chapter's Summit (\probref{c20-summit}), where a \emph{Grundy value} rather than a
boolean is rerooted -- see \thmref{c27-hackenbush} and \chref{c27} for the game
machinery this problem borrows the flavour of.
\end{topicintro}

\subsection{Rerooting a Count, and Reading It per Edge}

\begin{definition}{The count state and its per-edge verdict}{c20-count-state}
Assign each node a weight $w[v]\in\{0,1\}$ ($1$ = marked). Root the tree and set
the down fold $\mathit{down}[v] = \sum_{u\in\subtree(v)} w[u]$ -- a plain sum, so
$\mathit{TOTAL}=\mathit{down}[\mathit{root}]$ is the global count and the fold is
invertible. For an edge $(v,c)$ with $c$ on the child side, the two \emph{rerooted
counts} are the marks in each component of $T - \{(v,c)\}$:
\[
  \text{child side} = \mathit{down}[c], \qquad
  \text{parent side} = \mathit{TOTAL} - \mathit{down}[c].
\]
A per-edge predicate $P(c\!\to\! v)$ (``may move from $c$ onto $v$'') is a
decision on these two counts and the endpoint colours; the per-node verdict
$\mathit{ans}[v]\in\{0,1\}$ is then a reachability question on the directed graph
$\{(a,b): P(a\!\to\!b)\}$.
\end{definition}

\noindent Two-DFS did double duty in \secref{c20-twodfs}; here the up pass is so
cheap we do not even materialise an $\mathit{up}[\cdot]$ array -- the parent-side
count at every edge is just $\mathit{TOTAL}-\mathit{down}[c]$, one subtraction,
valid because the combine is a sum. The genuinely new step is the \emph{last}
line of \defref{c20-count-state}: once the counts are known on both sides of every
edge, the answers are no longer the counts but a Boolean functional of the
directed-edge graph they induce, so a multisource sweep -- not another fold --
finishes the job.

\begin{theorem}{Every-edge rerooted counts, then verdicts, cost $\bigO{n}$}{c20-count-edge-transfer}
For a count state as in \defref{c20-count-state}, all $2(n-1)$ directed-edge
counts $\mathit{down}[c]$ and $\mathit{TOTAL}-\mathit{down}[c]$ are available after
one post-order DFS, and assembling the per-node verdicts by a single
multisource reachability over the induced arcs costs $\bigO{n}$ in total.
\end{theorem}

\begin{proofbox}[\thmref{c20-count-edge-transfer}: one fold, one subtraction, one BFS]
The post-order fold computes $\mathit{down}[v]$ for all $v$ in $\bigO{n}$
(\thmref{c10-subtree-template}); $\mathit{TOTAL}=\mathit{down}[\mathit{root}]$.
For each edge $(v,c)$, whichever component is the child's is $\mathit{down}[c]$ and
the complement is $\mathit{TOTAL}-\mathit{down}[c]$ -- a subtraction, so $\bigO{1}$
per directed edge and $\bigO{n}$ over all edges; this \emph{is} the invertible
up-pass of \defref{c20-parent-piece} specialised to a sum, needing no sibling loop
and no prefix/suffix. Emitting each directed arc $a\!\to\!b$ is $\bigO{1}$ from
those two counts and the endpoint marks, so the arc set is built in $\bigO{n}$.
Finally, a multisource BFS/DFS from the marked nodes over that arc graph touches
each arc once, $\bigO{n}$, and marks exactly the nodes that reach a mark -- the
verdicts. All phases are $\bigO{n}$, so the total is $\bigO{n}$. \qed
\end{proofbox}

\begin{insightbox}[Why a ``game'' problem is a count problem]
CF~1626E's rules look combinatorial: hop toward a black vertex, and you may not
pick the same black two moves running. The whole difficulty is the alternation --
but it collapses to a counting test. You can safely step $x\!\to\!y$ (so that your
\emph{next} move is not forced back along the same edge) exactly when the far
side $y$'s component holds at least two black vertices, or $y$ is itself black.
That is a statement about \emph{how many} blacks sit past the edge, i.e.\ about
$\mathit{TOTAL}-\mathit{down}$, not about game trees. Rerooting the count turns the
rule into a per-edge yes/no, and reachability turns those into per-node answers.
This is the same ``count first, then read a verdict'' shape that Ch.~14 used for
its supply/demand DP (\thmref{c27-amortize}); the reroot supplies the
\emph{every}-edge counts that a single-root DFS would not give you for the
parent side.
\end{insightbox}

\subsection{The Algorithm: Count Reroot, Then Multisource Reachability}

\begin{lstlisting}[caption={Reroot a count of marked nodes: the parent side of every edge is TOTAL - down[c] (one subtraction), the counts drive a per-edge decision, and a multisource sweep assembles per-node verdicts.}, label={lst:c20-count}]
// Combine is a plain SUM -> invertible -> the up value is TOTAL - down[c].
// down[v] = count of marks in subtree(v)      (post-order, Ch.13 fold).
// For edge (v, c), c child of v:
//     marks past c   = down[c]
//     marks past v   = TOTAL - down[c]
// A per-node verdict is NOT the count: it is reachability on the directed
// arcs each edge test admits.
vector<vector<int>> adj;
vector<int> down_, w, par_;
int TOTAL;

void dfs_down(int v, int p) {                // post-order: subtree mark count
    par_[v] = p;
    down_[v] = w[v];
    for (int c : adj[v]) if (c != p) { dfs_down(c, v); down_[v] += down_[c]; }
}
// driver:
//   dfs_down(root, -1);  TOTAL = down_[root];
//   for each edge (v, c):
//       childSide = down_[c];  parentSide = TOTAL - down_[c];
//       emit arcs a->b from {childSide, parentSide, w[a], w[b]};
//   multisource-BFS/DFS over emitted arcs from the marked nodes -> ans[].
// O(n) total: one fold + O(1) subtraction per edge + one arc sweep.
// This skeleton is instantiated verbatim by lst:c20-1626e-soln.
\end{lstlisting}

\begin{figure}[H]
\centering
\begin{tikzpicture}[diagram, scale=0.92, transform shape,
    side/.style={circle, draw=#1, line width=0.9pt, fill=#1!10,
                 minimum size=8mm, font=\small\bfseries, text=#1!60!black},
    blk/.style={circle, draw=warningred, fill=warningred!85, line width=0.9pt,
                 minimum size=3.2mm, inner sep=0pt}]
  % the cut edge
  \node[side=primaryblue] (v) at (0,0)  {$v$};
  \node[side=successgreen] (c) at (3.1,0) {$c$};
  \draw[tedge, line width=1.1pt] (v) -- (2.05,0);
  \draw[tedge, line width=1.1pt] (2.55,0) -- (c);
  \node[annot, font=\scriptsize, text=warningred!70!black] at (2.3,0.45) {cut};
  % child side blob
  \node[blk] (m1) at (4.4,-0.6) {};
  \node[annot, font=\scriptsize, anchor=west] at (4.7,-0.6) {$\mathit{down}[c]$ marks past $c$};
  \draw[dashed, successgreen!60!black, line width=0.8pt, rounded corners]
       (3.35,-1.05) rectangle (7.6,0.55);
  % parent side blob
  \node[blk] (m2) at (-1.3,-0.6) {};
  \node[blk] (m3) at (-2.1,-0.2) {};
  \node[annot, font=\scriptsize, anchor=east] at (-1.0,-0.6) {$\mathit{TOTAL}-\mathit{down}[c]$ marks past $v$};
  \draw[dashed, primaryblue!60!black, line width=0.8pt, rounded corners]
       (-2.5,-1.05) rectangle (-0.35,0.55);
  \node[annot, font=\scriptsize, text=gray] at (2.3,-1.55)
       {$\text{child side}+\text{parent side}=\mathit{TOTAL}$ (a sum: the split is exact)};
\end{tikzpicture}
\caption{Cutting edge $(v,c)$ splits the marks into two rerooted counts,
$\mathit{down}[c]$ (the child subtree, unchanged by re-rooting) and
$\mathit{TOTAL}-\mathit{down}[c]$ (the parent side -- the up value, free because
the combine is a sum). Both numbers land on the same edge, so the two-sided
count that CF~1626E's move rule needs is available for \emph{every} edge in one
post-order pass.}\label{fig:c20-countfig}
\end{figure}

\subsection{Worked Example: Black $=\{4,10\}$ on the Running Tree $T$}

\begin{example}{Rerooting the mark count on $T$ and reading the verdicts}{c20-count-run}
Reuse $T$ (root $1$, $n=10$, edges $1\text{--}2,1\text{--}3,2\text{--}4,2\text{--}5,
3\text{--}6,3\text{--}7,5\text{--}8,7\text{--}9,9\text{--}10$). Declare vertices
$4$ and $10$ black ($\mathit{TOTAL}=2$), all else white. Post-order gives
$\mathit{down}[v]=$ the number of blacks in $\subtree(v)$; the parent side of edge
$(v,c)$ is $2-\mathit{down}[c]$. The safe-move arcs $a\!\to\!b$ are exactly those
where $b$'s far side holds $\ge 2$ blacks or $b$ is black; a multisource sweep from
$\{4,10\}$ yields the verdict per node. Every value below is hand-computed.
\end{example}

\begin{center}
\footnotesize
\renewcommand{\arraystretch}{1.2}
\begin{tabular}{l|cccccccccc}
\toprule
node $v$ & 1 & 2 & 3 & 4 & 5 & 6 & 7 & 8 & 9 & 10 \\
\midrule
$w[v]$ (black?) & 0 & 0 & 0 & 1 & 0 & 0 & 0 & 0 & 0 & 1 \\
$\mathit{down}[v]$ & 2 & 1 & 1 & 1 & 0 & 0 & 1 & 0 & 1 & 1 \\
arc destination $v$, sources into $v$ & -- & $\{5\}$ & $\{6\}$ & $\{2\}$ & $\{8\}$ & $\{1\}$ & -- & -- & -- & $\{9\}$ \\
$\mathit{ans}[v]$ (reaches black) & 0 & 1 & 0 & 1 & 1 & 0 & 0 & 1 & 1 & 1 \\
\bottomrule
\end{tabular}
\end{center}

\noindent The verdicts on $T$ are therefore $(0,1,0,1,1,0,0,1,1,1)$: nodes $4,10$
are black (verdict trivially $1$), nodes $2,5,8$ drain into $4$ and nodes $9$ into
$10$, while the interior nodes $1,3,6,7$ are stranded (they oscillate, never
closing in on a black).

\begin{tracebox}[Count reroot and sweep on $T$ (blacks $4,10$)]
Post-order $\mathit{down}[v]=\sum_{u\in\subtree(v)}w[u]$: leaves $4,10$ are black
so $\mathit{down}[4]{=}\mathit{down}[10]{=}1$; then $\mathit{down}[8]{=}0$,
$\mathit{down}[5]{=}w5{+}\mathit{down}[8]{=}0$, $\mathit{down}[2]{=}w4{+}\mathit{down}
[5]{=}1$; $\mathit{down}[9]{=}1$, $\mathit{down}[7]{=}\mathit{down}[9]{=}1$,
$\mathit{down}[6]{=}0$, $\mathit{down}[3]{=}\mathit{down}[6]{+}\mathit{down}[7]{=}1$;
$\mathit{down}[1]{=}\mathit{down}[2]{+}\mathit{down}[3]{=}2=\mathit{TOTAL}$. \checkmark\\
Edge cut counts (child side / parent side), e.g.\ $(2,5)$: $0/2$ -- past $5$ there
are $0$ blacks but past $2$ there are $2$, so a chip at $5$ may step to $2$'s side
(arc $5\!\to\!2$). $(5,8)$: $0/2$ gives arc $8\!\to\!5$. $(9,10)$ and $(7,9)$: $1/1$
so only $9\!\to\!10$ (destination $10$ is black) is admitted; $7\!\to\!9$ is not
($9$'s far side has $1$ black, not $\ge2$).\\
Sweep from $\{4,10\}$ over arcs into a reached node: $10$ admits $9$, $9$'s set is
empty; $4$ admits $2$, then $2$ admits $5$, $5$ admits $8$. Reached
$=\{4,10,2,9,5,8\}$ $\Rightarrow$ $\mathit{ans}=(0,1,0,1,1,0,0,1,1,1)$. \checkmark
\end{tracebox}

\subsection{Variations and Twists}

\begin{itemize}[nosep]
  \item \textbf{Count to verdict:} the same skeleton answers ``for every node,
        does some marked node lie within $k$ edges?'' -- swap the $\ge2$ test for
        a bounded-distance reroot (\thmref{c20-edge-transfer}). \hfill\lv{4}
  \item \textbf{Many marks, threshold $t$:} keep $\mathit{down}$ / $\mathit{TOTAL}
        -\mathit{down}$ and compare against $t$ rather than $2$; the reachability
        step is unchanged. \hfill\lv{4}--\lv{5}
  \item \textbf{Verdict needs a direction:} \thmref{c20-count-edge-transfer}'s
        arcs are asymmetric -- $a\!\to\!b$ may hold while $b\!\to\!a$ fails, so
        never assume bidirectionality when building the arc set. \hfill\lv{5}
  \item \textbf{Game value instead of boolean:} reroot a \emph{Grundy number}
        rather than a mark count -- that is the Summit (\probref{c20-summit}) and
        Ch.~14's \thmref{c27-hackenbush}, where the per-node object is richer
        than $\{0,1\}$. \hfill\lv{6}
\end{itemize}

\subsection{Common Pitfalls}

\begin{warningbox}[Trap 1 -- assuming the move test is symmetric]
The safe-cross predicate lives on the \emph{destination}'s far-side count, so
$a\!\to\!b$ and $b\!\to\!a$ use different counts and can disagree. Build both
directed arcs independently; a bidirectional (undirected) reachability sweep
over-counts answers and gives $1$ for stranded nodes like $1,3,6,7$ in
\exref{c20-count-run}.
\end{warningbox}
\begin{warningbox}[Trap 2 -- $\mathit{TOTAL}$ from the wrong root]
$\mathit{TOTAL}=\mathit{down}[\text{root}]$ only after a full post-order; reading a
partial $\mathit{down}$ and subtracting yields a negative parent-side count and a
garbage arc test. Compute every $\mathit{down}[v]$ first, \emph{then} form
$\mathit{TOTAL}-\mathit{down}[c]$.
\end{warningbox}
\begin{warningbox}[Trap 3 -- BFS on the wrong graph]
The verdicts come from a sweep over the \emph{transposed} arcs (from each node,
walk to the sources that may step \emph{onto} it), started at the black nodes.
Running the sweep on the forward arcs instead asks ``who can this node reach''
rather than ``who reaches this node'' -- the reachability direction is exactly the
move direction, so transposing (or equivalently starting from targets) is the
correct, single-source-free formulation.
\end{warningbox}

\subsection{Problem Lab}

\solvedproblem{c20-count-solved}{Black and White Tree}{Codeforces 1626E}
  {https://codeforces.com/problemset/problem/1626/E}{CF 2400 · L5}
  {count reroot (subtree marks and TOTAL-minus-subtree), per-edge two-sided
   counts, directed move graph, transposed multisource reachability, game flavour}

\begin{problemstatement}[Black and White Tree (statement and sample verified against the live problem)]
A tree of $n$ vertices ($3 \le n \le 3\cdot10^{5}$); some vertices ($\ge 2$) are
black ($c_i=1$), the rest white. Place a chip at vertex $i$. A move: from the
current vertex $x$, choose a black vertex $y$ and move the chip along the first
edge of the $x\!\to\!y$ path. You may \emph{not} choose the same black $y$ on two
consecutive moves. The chip is saved once it lands on a black vertex (a chip
starting on a black vertex needs no moves), and the process is abandoned if it
exceeds $100^{500}$ moves. For every vertex $i$, print $1$ if some legal sequence
of moves brings the chip to a black vertex, else $0$.\\
\textbf{Sample}: $n=8$, colours $(0,1,0,0,0,0,1,0)$ (vertices $2,7$ black), edges
$8\text{--}6,\,2\text{--}5,\,7\text{--}8,\,6\text{--}5,\,4\text{--}5,\,6\text{--}1,\,
7\text{--}3$; the output is \texttt{0 1 1 1 1 0 1 1} (vertices $1$ and $6$ fail).
\end{problemstatement}

\begin{attemptgate}[Attempt it yourself first -- \lv{5}, give yourself 45 minutes]
The alternation rule is the only obstacle; forget ``game tree'' and ask a counting
question. (a)~From $x$, if you step toward $y$ onto neighbour $b$, when are you
\emph{not} immediately forced straight back? (b)~That depends only on how many
black vertices lie past the edge $x\text{--}b$, on $b$'s side. Count them.
\peek{Root the tree and let $\mathit{down}[v]$ be the black count in
      $\subtree(v)$. For a neighbour $b$ of $x$, what is the black count on $b$'s
      side of the cut edge?}
\nudge{You can move $x\!\to\!b$ and not be forced back iff $b$ is black, or the
      component holding $b$ (after deleting edge $x\text{--}b$) has $\ge 2$ black
      vertices -- one to make the move, one to continue. Treat each such legal move
      as a directed arc $x\!\to\!b$.}
\fullhint{Build arcs $x\!\to\!b$ for every legal move (destination $b$: $c_b=1$ or
the far-side count $>1$, where far side is $\mathit{down}[b]$ if $b$ is a child of
$x$ else $\mathit{TOTAL}-\mathit{down}[x]$). A vertex is a winner iff it can reach
\emph{some} black vertex along these arcs. Equivalently, transpose the arc graph
(reverse every arc) and run a \emph{multisource} BFS/DFS whose sources are all the
black vertices -- a node is marked reachable exactly when a black vertex is
reachable from it in the original. Output reachability. Everything is $\bigO{n}$:
one post-order for $\mathit{down}$, one arc-build sweep, one BFS.}
\end{attemptgate}

\solstep{Understanding the Problem}
Two facts make this a counting problem, not a game-tree problem. First, the chip's
only real constraint is the ``different black each move'' rule; a black vertex is
merely a \emph{target you alternate between}, so what matters locally is whether,
at your current vertex, there are \emph{enough} black vertices off in a chosen
direction to keep switching. Second, once the chip crosses an edge it can never be
pushed back by the rule unless that side has too few blacks. So a move
$x\!\to\!b$ is \emph{safe} (you can land on $b$ and still make a legal next move
that does not undo the step) precisely when $b$ is black or $b$'s side of the
$x\!\text{--}\,b$ cut contains at least two black vertices. Counting blacks per
component is a subtree count -- and the complement of a subtree is the parent-side
rerooted count $\mathit{TOTAL}-\mathit{down}[\cdot]$. That count answers the safety
question for \emph{every} directed edge in one pass, and the ``can reach a black''
answer per node is just reachability in the safe-move digraph.

\solstep{Key Observations}
\begin{observation}
A chip at any vertex can always move \emph{toward} some black vertex; the only
failure mode is being forced back along the edge it just crossed. So legality
reduces to, per edge, whether the far side offers ``two distinct black targets''
or the endpoint is already black.
\end{observation}
\begin{observation}
The number of black vertices on $b$'s side of the edge $(x,b)$ is a rerooted
count: $\mathit{down}[b]$ if $b$ is a child of $x$, and $\mathit{TOTAL}-
\mathit{down}[x]$ if $b$ is $x$'s parent. Because the mark combine is a sum, the
parent side costs one subtraction -- no up array, no prefix/suffix
(\defref{c20-count-state}).
\end{observation}
\begin{observation}
Vertex $v$ is a winner ($\mathit{ans}[v]=1$) iff $v$ reaches a black vertex along
the safe-move arcs. Transposing the arcs and running a \emph{multisource} BFS from
the black vertices computes all winners at once -- the two cases $x\!\to\!y$ and
$y\!\to\!x$ differ, so the graph is genuinely directed.
\end{observation}

\solstep{Think It Through With Me}
\begin{thinkbox}
\deadend{Simulate the game: from each vertex, alternate between the two nearest
        black vertices, and detect a draw with a state graph of size $n\cdot B$.}
{The state must remember the previously chosen black, so it is $\bigO{n B}$ -- up to
$\bigO{n^2}$, and $n=3\cdot10^5$ makes it hopeless (\tle). It also misses that the
answer only depends on \emph{counts past each edge}, not on which specific black
was last used.}
\deadend{Reroot the subtree black counts and stop, reading $\mathit{ans}[v]$ off
        $\mathit{down}[v]$ directly.}
{$\mathit{down}[v]$ counts blacks \emph{below} $v$, but a winner can be decided by
blacks \emph{above} it; the count is per-edge and per-direction, and the verdict is
a reachability question over those directed edges, not a single-node value. The up
count $\mathit{TOTAL}-\mathit{down}$ is only half -- you still sweep the arcs.}
\aha{Turn each safe move into a directed arc (destination's far side has $>1$ black
or is black), then transpose and multisource-BFS from the black vertices. The
subtree count plus its subtraction gives both directions' counts on every edge in
$\bigO{n}$; the sweep gives all verdicts in $\bigO{n}$. That is
\coderef{c20-count} instantiated by \coderef{c20-1626e-soln}.}
\end{thinkbox}

\solstep{Solution}
Root at $1$; post-order computes $\mathit{down}[v]$ (blacks in $\subtree(v)$) and
the parent array; $\mathit{TOTAL}=\mathit{down}[1]$. For every ordered adjacent
pair $(i,j)$ decide the transposed arc: $j$ may move onto $i$ (add $j$ to
$\mathit{arc}[i]$) iff $c_i=1$ or the black count on $i$'s side of edge $(i,j)$
exceeds $1$ -- that side is $\mathit{TOTAL}-\mathit{down}[j]$ when $j$ is a child of
$i$ and $\mathit{down}[i]$ when $j$ is $i$'s parent. Finally run a multisource
BFS over $\mathit{arc}$ seeded with every black vertex; a node is reached iff it
can, in the original graph, move onto a chain ending at a black vertex. Print the
reachability bit per node.

\solstep{Proof of Correctness}
\begin{proofbox}[\coderef{c20-1626e-soln} prints $1$ exactly for the winning vertices]
\emph{Arc characterisation.} A chip moving $x\!\to\!b$ is forced straight back off
$b$ iff, after landing, it must target a black vertex whose path's first edge is
$b\!\to\!x$; with the ``no repeat'' rule and $x$'s last target being the black it
chose, this is unavoidable exactly when $b$'s component of $T-\{(x,b)\}$ has at
most one black and $b$ is not itself black. Hence the move is safe iff $b$ is black
or $b$'s side has $\ge 2$ blacks -- and $b$'s side count is, by
\defref{c20-count-state}, $\mathit{down}[b]$ (child) or $\mathit{TOTAL}-\mathit{down}
[x]$ (parent), both computed correctly in the post-order. So $\mathit{arc}$ (built
from this test with $i$ as destination) records exactly the safe moves.
\emph{Reachability $=$ winning.} A vertex wins iff some safe-move path reaches a
black vertex: forward, the chip follows those arcs and lands on a black; backward,
any winning sequence is a chain of safe moves (each step, by definition, was not
forced back), so it is a directed path to a black. Marking the black vertices and
traversing the \emph{transposed} arcs reaches precisely the vertices that have a
directed path to a black in the original graph. Thus $\mathit{used}[v]=1$ iff $v$
is a winner. \qed
\end{proofbox}

\solstep{Implementation}
\begin{lstlisting}[caption={Codeforces 1626E -- reroot the black count (down[c] / TOTAL-minus), turn each safe move into a directed arc, transposed multisource BFS from the blacks; O(n)}, label={lst:c20-1626e-soln}]
#include <bits/stdc++.h>
using namespace std;
int n;
vector<int> col, cntb, par;
vector<vector<int>> adj, arc;   // arc[i]: nodes that may move ONTO i
vector<char> used;

void dfs(int v, int p) {                    // post-order black count
    par[v] = p;
    cntb[v] = col[v];
    for (int c : adj[v]) if (c != p) { dfs(c, v); cntb[v] += cntb[c]; }
}
int main() {
    scanf("%d", &n);
    adj.assign(n + 1, {}); arc.assign(n + 1, {});
    col.assign(n + 1, 0); par.assign(n + 1, 0);
    cntb.assign(n + 1, 0); used.assign(n + 1, 0);
    for (int i = 1; i <= n; i++) scanf("%d", &col[i]);
    for (int i = 0; i < n - 1; i++) {
        int u, v; scanf("%d %d", &u, &v);
        adj[u].push_back(v); adj[v].push_back(u);
    }
    dfs(1, 0);
    int total = cntb[1];
    for (int i = 1; i <= n; i++)
        for (int j : adj[i]) {
            int iSide;                       // blacks on i's side of edge (i,j)
            if (par[j] == i) iSide = total - cntb[j];   // j is child of i
            else             iSide = cntb[i];           // j is parent of i
            if (col[i] == 1 || iSide > 1) arc[i].push_back(j);  // j may step to i
        }
    queue<int> q;                            // multisource BFS from every black
    for (int i = 1; i <= n; i++) if (col[i] == 1) { used[i] = 1; q.push(i); }
    while (!q.empty()) {
        int v = q.front(); q.pop();
        for (int j : arc[v]) if (!used[j]) { used[j] = 1; q.push(j); }
    }
    for (int i = 1; i <= n; i++) printf("%d ", (int)used[i]);
    puts("");
    return 0;
}
\end{lstlisting}

\solstep{Trace}
\begin{tracebox}[Official sample: $n=8$, blacks $2,7$, edges as stated $\Rightarrow$ \texttt{0 1 1 1 1 0 1 1}]
Post-order from root $1$: children $1\!\to\!6\!\to\!\{8,5\}$, $8\!\to\!7\!\to\!3$,
$5\!\to\!\{2,4\}$. Black counts: $\mathit{down}[2]{=}1$, $\mathit{down}[7]{=}1$,
$\mathit{down}[4]{=}\mathit{down}[3]{=}\mathit{down}[8_{\text{self}}]{=}0$,
$\mathit{down}[5]{=}1$, $\mathit{down}[8]{=}1$, $\mathit{down}[6]{=}2$,
$\mathit{down}[1]{=}2=\mathit{total}$.\\
Arcs (destination $\Leftarrow$ sources): onto black $2$ comes from $5$; onto black
$7$ from $8$ and $3$; onto $5$ from $4$ (its far side past $5$ has $2$ blacks); onto
$6$ from $1$ (far side past $6$ has $2$). Stranded interiors: edge $(6,5)$ splits
$1/1$, so neither $5\!\to\!6$ nor $6\!\to\!5$ is safe; edge $(8,6)$ splits $1/1$
likewise.\\
BFS from $\{2,7\}$: $2$ marks $5$; $7$ marks $8,3$; $5$ marks $4$; done. Reached
$\{2,3,4,5,7,8\}$, unreached $\{1,6\}$ (nothing safely reaches them). Output
$\boxed{0\,1\,1\,1\,1\,0\,1\,1}$ \checkmark\ matches CF's sample.
\end{tracebox}

\solstep{Complexity}
\begin{complexitybox}
Time $\bigO{n}$: one post-order DFS for $\mathit{down}$, one sweep over the
$2(n-1)$ directed edges to build $\mathit{arc}$ (each an $\bigO{1}$ count test), and
one multisource BFS/DFS touching each arc once. Memory $\bigO{n}$ for the
adjacency, arc lists and the three integer arrays. $\mathit{down}$ fits in $\mathit
{int}$ (a count $\le n$); no overflow, no modulo. Recursion depth $h$ -- at
$n\le3\cdot10^5$ convert the post-order to the explicit stack of
\coderef{c09-iter} to avoid a stack overflow. Verdict: \ac; the count-reroot
template -- the up value is free as $\mathit{TOTAL}-\mathit{down}$, and the
interesting work moves to the directed reachability that reads the verdicts.
\end{complexitybox}

\solstep{Pitfalls and Variants}
\begin{warningbox}
Ways to lose CF~1626E: (i)~making the move graph \emph{undirected} -- the arc test
is asymmetric, so build both directions separately. (ii)~using $\mathit{down}[c]$
for the parent side (the far side past $v$ is $\mathit{TOTAL}-\mathit{down}[c]$, a
different number). (iii)~forgetting that a black destination is always enterable
($\mathit{col}[i]=1$ short-circuits the count test) and rejecting chips that sit one
edge from a black. (iv)~deep recursion on a path tree -- iterate the post-order.
Variants: replace ``$\ge2$ blacks'' with any per-edge predicate on the rerooted
count to solve ``for every node, can it reach a node with property $P$'' families;
swap the black \emph{count} for a \emph{Grundy value} and you are exactly in the
Summit's territory (\probref{c20-summit}, \thmref{c27-hackenbush}), where the
per-node object is a game number rather than a bit.
\end{warningbox}

\subsection{Practice Problems}

\begin{stuckbox}
\begin{itemize}[nosep, leftmargin=1.4em]
  \item Did you compute \emph{both} the subtree count $\mathit{down}[c]$ and the
        complement $\mathit{TOTAL}-\mathit{down}[c]$ for every directed edge?
  \item Is the move graph built \emph{directed} (destination's far-side count
        drives the test), and did you seed the BFS from all black vertices at once?
  \item Did you short-circuit ``destination is black'' before the count test?
\end{itemize}
\end{stuckbox}

\begin{practiceset}
  \practice{Labeling the Tree with Distances}{Codeforces 1794E}
    {https://codeforces.com/problemset/problem/1794/E}{CF 2400 · L5}
    {Given a tree and a target distance $d$, find every vertex $r$ that can serve
    as a root from which a chosen set of markers sits at the prescribed pairwise
    distances -- a per-root \emph{feasibility verdict}, decided by rerooting a
    distance/height count so that each node sees both its own subtree statistic
    ($\mathit{down}$) and the parent side ($\mathit{TOTAL}-\mathit{down}$, or the
    second-best height) and answers yes/no, exactly the shape of
    \defref{c20-count-state}. Samples (verified): $n=6$ report $2$ vertices
    ($2,4$); $n=5$ report $1$ vertex ($3$); $n=3$ report $0$.
    }
\end{practiceset}

\paragraph{Back-cites (URLs live elsewhere in the book).}
\begin{itemize}[nosep,leftmargin=1.4em]
  \item \textbf{CF~1324F ``Maximum White Subtree''} (\chref{c10} owns the URL): the
        count/weight reroot whose up value is a subtraction on the additive side
        and a $\max$ elsewhere -- reread it as a per-edge count like
        \defref{c20-count-state}.
  \item \textbf{Ch.~14 supply/demand reroot} (\chref{c27} owns those URLs): the
        ``per-edge, count-driven verdict'' pattern this section shares, before the
        game-value version reached in the Summit (\probref{c20-summit}).
\end{itemize}

\begin{recallbox}[10-second refresher -- 15.4 Rerooting a count]
Reroot a mark count: $\mathit{down}[c]$ is the subtree count and the parent side is
$\mathit{TOTAL}-\mathit{down}[c]$ -- a subtraction, since the combine is a sum, so
no up array and no prefix/suffix (\thmref{c20-count-edge-transfer}). Then read a
\emph{verdict}, not the count: build a \emph{directed} safe-move graph from the two
sides' counts (destination black, or far side $\ge 2$) and multisource-BFS from the
marked nodes to mark winners. CF~1626E (\probref{c20-count-solved}): on $T$ with
blacks $\{4,10\}$ the verdicts are $(0,1,0,1,1,0,0,1,1,1)$; the official $n=8$
sample gives \texttt{0 1 1 1 1 0 1 1}.
\selftest{
  \item Why is the parent-side mark count just $\mathit{TOTAL}-\mathit{down}[c]$?
  \item What makes the move graph directed rather than undirected?
  \item State the exact test for a safe move $x\!\to\!b$.
  \item Why seed the BFS from \emph{all} black vertices and traverse transposed arcs?
  \item Which node in \exref{c20-count-run} has verdict $0$ despite having a black
        in its own subtree, and why?
}
\end{recallbox}
